\section{Exponential equations for an individual density}
\label{sec:individual-exponential-equations}

Fix a nonzero nonnegative $f\in L^1(\mathrm{d}V)$. An exponential subsolution is a common lower bound for the continuous solutions of the truncated equations. Those solutions decrease, so \cref{thm:exponential-closedness} gives an exact equation at the limit. The limiting solution dominates every subsolution. A positive threshold $b_h(f)$ supplies the required lower bound for each exponent.

\begin{theorem}\label{thm:individual-exponential-subsolution}\index{Monge--Ampere equation@Monge--Amp\`ere equation!greatest exponential solution}
Fix $\lambda>0$. Suppose there is $s\in\PSH(X,h)$ such that
\[
 \mathrm{MA}^{\mathrm{ext}}_h(s)\geq\mathrm{e}^{\lambda s}f\,\mathrm{d}V.
\]
Then the equation
\[
 \mathrm{MA}^{\mathrm{ext}}_h(u_\lambda)
       =\mathrm{e}^{\lambda u_\lambda}f\,\mathrm{d}V
\]
has a greatest solution. More precisely, $u_\lambda$ is the greatest genuine $h$-plurisubharmonic subsolution of this equation, and its Monge--Amp\`ere measure is finite. It is the decreasing limit of the continuous solutions for $f_k=\min(f,k)$.
\end{theorem}
\begin{proof}
Let $u_k$ solve the continuous exponential equation for $f_k$, using \cref{thm:nguyen-exp}. By \cref{lem:exp-order}, $u_{k+1}\leq u_k$.

Every genuine subsolution $v$ for $f$ lies below each $u_k$. Indeed, for $t>0$, on $\{u_k<v-t\}$ we have
\[
 \mathrm{MA}_h(u_k)=\mathrm{e}^{\lambda u_k}f_k\,\mathrm{d}V
 \leq\mathrm{e}^{-\lambda t}\mathrm{e}^{\lambda v}f\,\mathrm{d}V
 \leq\mathrm{e}^{-\lambda t}\mathrm{MA}^{\mathrm{ext}}_h(v-t).
\]
The first potential $u_k$ is bounded, so $P_{h/2}(u_k)$ is genuine and $v-t\preceq u_k$. Strict domination, \cref{thm:singular-strict-domination}, gives $u_k\geq v-t$. Let $t\downarrow0$.

In particular, $s\leq u_k\leq u_1$. Thus the decreasing limit $u_\lambda$ is genuine. By \cref{lem:monotone-psh-capacity}, the convergence holds in capacity. Choose a common constant $A\geq\sup_Xu_1$. Apply \cref{thm:exponential-closedness} to $u_k-A$, with densities $\mathrm{e}^{\lambda A}f_k\to\mathrm{e}^{\lambda A}f$ in $L^1$. It proves the equation and finiteness of its left side. Finally every subsolution $v$ lies below all $u_k$, hence below $u_\lambda$, which is itself a subsolution.
\end{proof}

Every subsolution lies below this solution, so its extremal property is independent of the truncation used to construct it. Smaller solutions with other singularities are not excluded. For $f\in L^p$, $p>1$, the continuous solution gives a bounded lower barrier. The greatest solution is then bounded and agrees with it by bounded exponential comparison.

\Needspace{20\baselineskip}
\begin{corollary}\label{cor:individual-positive-threshold-exponential}
If $b=b_h(f)>0$, the preceding greatest solution exists for every $\lambda>0$. If $v\in\PSH(X,h)$ has $\sup_Xv=0$ and
$\mathrm{MA}^{\mathrm{ext}}_h(v)\geq bf\,\mathrm{d}V$, then
\[
 u_\lambda\geq v+\lambda^{-1}\log b.
\]
Writing
\[
 M_\lambda=\sup_Xu_\lambda,\qquad
 v_\lambda=u_\lambda-M_\lambda,\qquad
 c_\lambda=\mathrm{e}^{\lambda M_\lambda},
\]
we have
\[
 c_\lambda\longrightarrow b_h(f),\qquad
 \mathrm{e}^{\lambda u_\lambda}f\longrightarrow b_h(f)f
           \quad\text{in }L^1(\mathrm{d}V)
 \quad(\lambda\downarrow0).
\]
\end{corollary}
\begin{proof}
The attained-threshold theorem supplies $v$. Since $v\leq0$, the function $s=v+\lambda^{-1}\log b$ satisfies
\[
 \mathrm{MA}^{\mathrm{ext}}_h(s)\geq bf\,\mathrm{d}V
       \geq\mathrm{e}^{\lambda s}f\,\mathrm{d}V.
\]
The theorem gives $u_\lambda\geq s$, hence $c_\lambda\geq b$.

Set $g_\lambda=\mathrm{e}^{\lambda v_\lambda}f$. By \cref{lem:normalized-exponential-npp}, applied to the fixed finite measure $f\,\mathrm{d}V$, we have $g_\lambda\to f$ in $L^1$. Each $g_\lambda$ is nonzero, since a genuine quasi-plurisubharmonic function is finite almost everywhere. The equation for $u_\lambda$ implies
\[
 c_\lambda\leq b_h(g_\lambda)
\]
by the maximal subsolution characterization. Upper semicontinuity therefore gives
\[
 b\leq\liminf_{\lambda\downarrow0}c_\lambda
 \leq\limsup_{\lambda\downarrow0}c_\lambda
 \leq b_h(f)=b.
\]
This proves the scalar limit. Since $g_\lambda\to f$ in $L^1$, it also proves $c_\lambda g_\lambda\to bf$ in $L^1$.
\end{proof}

The right-hand-side measures therefore converge in total variation. This answers the scalar and measure questions in \cref{sec:zero-exponent-topology}. Under positive lower mass, \cref{cor:positive-lower-mass-exact-threshold} also gives capacity convergence along every $L^1$-convergent normalized subsequence and identifies an exact solution at $b_h(f)$. For $L^p$ data, $p>1$, \cref{prop:symmetric-equation-clusters} gives uniform precompactness. In either case, convergence of the whole normalized family requires identification of all cluster limits. Smooth strictly positive dual attainment gives whole-family uniform convergence in \cref{cor:smooth-dual-exponential-limit}. If the smooth maximizing form is only semipositive, but strictly positive outside a proper analytic set, \cref{prop:semipositive-dual-exponential-limit} gives whole-family $L^1$ convergence and a fixed lower barrier. For $L^p$ data, $p>1$, \cref{cor:semipositive-dual-lp-limit} strengthens this to uniform convergence.

\Needspace{7\baselineskip}
For a general positive individual threshold, the corollary gives
\[
 v_\lambda\geq v-\lambda^{-1}\log(c_\lambda/b_h(f)).
\]
The next estimate quantifies this scalar error when the density root is approximable by bounded roots. Whether an exponential subsolution exists when $b_h(f)=0$ is a separate question; \cref{thm:individual-exponential-subsolution} applies whenever such a subsolution is available.

\subsection{An error bound for the scalar limit}

Suppose that $f^{1/n}$ belongs to the space $\mathcal D_h^0$ of \cref{subsec:bounded-root-approximation}. The uniform tail condition of \cref{prop:root-uniform-tails} then gives a rate for the convergence $c_\lambda\to b_h(f)$ in \cref{cor:individual-positive-threshold-exponential}. Split the root $p=f^{1/n}$ at a height $K$. On $\{p\leq K\}$, the error $p(1-\mathrm{e}^{\lambda v/n})$ of the damped root has norm $O(K\lambda)$ in $\mathcal D_h$, uniformly over normalized potentials $v$, while on $\{p>K\}$ it is bounded by the tail of $p$, which is small in that norm. Optimizing over $K$ gives a modulus valid for every normalized potential.

\begin{proposition}\label{prop:exponential-root-tail-modulus}
Let $0\leq p\in\mathcal D_h^0$ be nonzero, and put $f=p^n$, $R=\rho_h(p)$ and $b=R^{-n}=b_h(f)$. There is a constant $C_h$, depending only on $h$ and $\mathrm{d}V$, such that the function
\[
 \eta_p(\lambda)=\inf_{K>0}
 \left(\|p\mathbf1_{\{p>K\}}\|_{\mathcal D_h}+C_hK\lambda\right)
\]
tends to zero as $\lambda\downarrow0$ and satisfies
\begin{equation}\label{eq:uniform-root-damping}
 \big\|p(1-\mathrm{e}^{\lambda v/n})\big\|_{\mathcal D_h}
 \leq\eta_p(\lambda)
 \qquad\bigl(v\in\PSH(X,h),\ \sup_Xv=0\bigr).
\end{equation}
Let $u_\lambda$, $v_\lambda$ and $c_\lambda$ be the greatest exponential solutions and their normalizations from \cref{cor:individual-positive-threshold-exponential}. If $\eta_p(\lambda)<R$, then
\[
 b\leq c_\lambda\leq\bigl(R-\eta_p(\lambda)\bigr)^{-n}
\]
and
\begin{align*}
 \big\|\mathrm{MA}^{\mathrm{ext}}_h(v_\lambda)
              -bf\,\mathrm{d}V\big\|_{\mathrm{TV}}
 &\leq(c_\lambda-b)\|p\|_{L^n}^{n}
       +nb\|p\|_{L^n}^{n-1}\eta_p(\lambda).
\end{align*}
If $p\in L^r$ for some finite $r>n$, the same bounds hold with $\eta_p(\lambda)$ replaced by $C_{h,r}\|p\|_{L^r}\lambda$, where $C_{h,r}$ depends only on $h$, $\mathrm{d}V$ and $r$. In particular, for each such $p$ the scalar and total-variation errors are $O(\lambda)$.
\end{proposition}
\Needspace{10\baselineskip}
\begin{proof}
By \cref{lem:uniform-exponential}, normalized $h$-plurisubharmonic functions are uniformly bounded in every $L^s$ with $s<+\infty$. By \eqref{eq:root-subcritical-bound} with $r=2n$, their norms $\|v\|_{\mathcal D_h}$ are therefore uniformly bounded; let $C_h$ be $1/n$ times this bound. Let $v$ be normalized and $K>0$. On $\{p\leq K\}$ we have
\[
 0\leq p(1-\mathrm{e}^{\lambda v/n})
     \leq K\lambda(-v)/n,
\]
since $1-\mathrm{e}^{-x}\leq x$ for $x\geq0$, while on $\{p>K\}$ we use the bound $p(1-\mathrm{e}^{\lambda v/n})\leq p$. The monotonicity of $\|\cdot\|_{\mathcal D_h}$ and the triangle inequality give
\[
 \big\|p(1-\mathrm{e}^{\lambda v/n})\big\|_{\mathcal D_h}
 \leq\|p\mathbf1_{\{p>K\}}\|_{\mathcal D_h}+C_hK\lambda,
\]
and taking the infimum over $K$ proves \eqref{eq:uniform-root-damping}. By \cref{prop:root-uniform-tails}, and since the $L^n$ tails of $p$ tend to zero, $\|p\mathbf1_{\{p>K\}}\|_{\mathcal D_h}\to0$ as $K\to\infty$. Choosing $K$ first and then letting $\lambda\downarrow0$ shows that $\eta_p(\lambda)\to0$.

Now set $q_\lambda=p\,\mathrm{e}^{\lambda v_\lambda/n}$, so that $q_\lambda^n=g_\lambda=\mathrm{e}^{\lambda v_\lambda}f$, with $g_\lambda$ as in the proof of \cref{cor:individual-positive-threshold-exponential}. Since $0\leq q_\lambda\leq p$ and $\rho_h(p-q_\lambda)\leq\eta_p(\lambda)$ by \eqref{eq:uniform-root-damping}, the triangle inequality and monotonicity give
\[
 R-\eta_p(\lambda)\leq\rho_h(q_\lambda)\leq R.
\]
As in the proof of \cref{cor:individual-positive-threshold-exponential}, $b\leq c_\lambda\leq b_h(g_\lambda)=\rho_h(q_\lambda)^{-n}$, which yields the scalar bound. For the measures, the factorization of $p^n-q_\lambda^n$ gives $0\leq p^n-q_\lambda^n\leq np^{n-1}(p-q_\lambda)$, and H\"older's inequality implies
\[
 \|f-g_\lambda\|_{L^1}
 \leq n\|p\|_{L^n}^{n-1}\|p-q_\lambda\|_{L^n}
 \leq n\|p\|_{L^n}^{n-1}\eta_p(\lambda).
\]
Since $\mathrm{MA}^{\mathrm{ext}}_h(v_\lambda)=c_\lambda g_\lambda\,\mathrm{d}V$, the identity $c_\lambda g_\lambda-bf=(c_\lambda-b)g_\lambda-b(f-g_\lambda)$, together with $\|g_\lambda\|_{L^1}\leq\|f\|_{L^1}=\|p\|_{L^n}^n$, proves the total-variation bound.

Finally, suppose $p\in L^r$ with $n<r<+\infty$, and choose $n<s<r$. By \eqref{eq:root-subcritical-bound} with exponent $s$, and since $\|\cdot\|_{L^n}\leq\|\cdot\|_{L^s}$, we have $\|\cdot\|_{\mathcal D_h}\leq(1+C_s)\|\cdot\|_{L^s}$. H\"older's inequality and \cref{lem:uniform-exponential} then give, for every normalized $v$,
\[
 \|p(-v)\|_{\mathcal D_h}
 \leq(1+C_s)\|p(-v)\|_{L^s}
 \leq(1+C_s)\|p\|_{L^r}\|v\|_{L^{rs/(r-s)}}
 \leq nC_{h,r}\|p\|_{L^r}.
\]
Since $0\leq1-\mathrm{e}^{\lambda v/n}\leq\lambda(-v)/n$, the left side of \eqref{eq:uniform-root-damping} is therefore at most $C_{h,r}\|p\|_{L^r}\lambda$, and the scalar and total-variation estimates can be repeated with this bound in place of $\eta_p(\lambda)$. For $\lambda$ so small that $C_{h,r}\|p\|_{L^r}\lambda\leq R/2$, the function $t\mapsto(R-t)^{-n}$ has bounded derivative on $[0,R/2]$, so the scalar error $c_\lambda-b$, and hence the total-variation error, is $O(\lambda)$.
\end{proof}

For a general $p\in\mathcal D_h^0$ the modulus $\eta_p(\lambda)$ tends to zero, but it need not be $O(\lambda)$. The rescaled scalar error $\lambda^{-1}\log(c_\lambda/b)$ is the loss in the lower bound $v_\lambda\geq v-\lambda^{-1}\log(c_\lambda/b)$ for normalized threshold subsolutions $v$, and the scalar bound of the proposition gives
\[
 0\leq\lambda^{-1}\log(c_\lambda/b)
 \leq-\frac n\lambda\log\bigl(1-\eta_p(\lambda)/R\bigr),
\]
which stays bounded as $\lambda\downarrow0$ when $\eta_p(\lambda)=O(\lambda)$, in particular for $p\in L^r$ with $r>n$; root tails alone do not give such a bound. Indeed, on the products of \cref{sec:zero-lower-mass-product-equations} there are densities whose roots lie in $\mathcal D_h^0$ and whose rescaled errors $\lambda^{-1}\log(c_\lambda/b_h(f))$ tend to $+\infty$ (\cref{sec:nonlp-selection}); for their roots $\eta_p(\lambda)/\lambda$ is unbounded. Even when a uniform lower barrier for $v_\lambda$ is available, passing to the limit in capacity requires the singularity-type or lower-mass hypotheses of a continuity theorem. The proposition estimates only scalars and measures; when the lower mass vanishes, it does not give an exact equation at the threshold.


\Needspace{12\baselineskip}
\subsection{Exponential equations at zero threshold}
\label{subsec:zero-threshold-product-exponential}

Let $(X_0,h_0)$ be a connected compact Hermitian manifold of dimension $m\geq1$ with $v_-(h_0)=0$. Put $Y=\mathbf P^k$, $k\geq1$, with normalized Fubini--Study form $\kappa$, and set $X=X_0\times Y$, $H=h_0+\kappa$, $N=m+k$ and
\[
 C=\binom Nm,\qquad V_0=\int_{X_0}h_0^m,\qquad
 \mu=V_0^{-1}h_0^m\wedge\kappa^k,\qquad R=\mu/\mathrm dV.
\]
Here $\mathrm dV$ is any smooth strictly positive probability volume form on $X$, and factor pullbacks are understood. The potential $u_\alpha=1-(1-\psi)^\alpha$ and its density $q_\alpha$ are those of \cref{lem:projective-sublogarithmic-density}, for fixed $0<\alpha<1$ and pole $o=[1:0:\cdots:0]$. Its exponential weight can absorb a concentration. Choosing the projective components accordingly gives one zero-threshold density with genuine exponential solutions for every positive exponent. On $Y\setminus\{o\}$ put $L=1-\psi\geq1$, $t=\mathrm e^{1-L}$ and
\[
 a=\alpha L^{\alpha-1},\qquad Q=t+(1-t)a,\qquad
 D=t(1-a)+(1-t)\alpha(1-\alpha)L^{\alpha-2},
\]
\[
 W=kQ^{k-1}D,\qquad
 B_\tau=q_\alpha\mathrm e^{\tau(L^\alpha-1)}\quad(\tau>0),\qquad
 b_0=(1-\alpha)^k.
\]

\begin{lemma}\label{lem:projective-exponential-barrier-growth}
The radial expressions are
\[
 q_\alpha=Q^{k-1}D\,\mathrm e^{k(L-1)},\qquad q_\alpha\geq b_0,
 \qquad W(L)\sim k\alpha^k(1-\alpha)L^{-1-k(1-\alpha)}.
\]
For every $\tau>0$, $\mathrm e^{\tau(L^\alpha-1)}W(L)\to+\infty$ as $L\to+\infty$. In particular, $B_\tau\kappa^k$ has infinite mass on every punctured neighbourhood of $o$.
\end{lemma}
\begin{proof}
In the radial calculation of \cref{lem:projective-sublogarithmic-density}, $t=r/(1+r)$ and $-dQ/dL=D$. The Fubini--Study radial probability is $k\mathrm e^{-k(L-1)}\mathrm dL$, whereas the radial mass of $q_\alpha\kappa^k$ is $W(L)\mathrm dL$. Their ratio gives the formula for $q_\alpha$. The lower bound follows from $\kappa+\ddc u_\alpha\geq(1-\alpha)\kappa$ away from the pole. Since $Q\sim\alpha L^{\alpha-1}$ and $D\sim\alpha(1-\alpha)L^{\alpha-2}$, the displayed asymptotic follows. The logarithm of the weighted radial density is $\tau L^\alpha-[1+k(1-\alpha)]\log L+O(1)$, which tends to infinity. Its integral over every tail is therefore infinite.
\end{proof}

\begin{theorem}\label{thm:product-zero-threshold-exponential-criterion}\index{Monge--Ampere equation@Monge--Amp\`ere equation!zero threshold criterion}
Let $f\geq0$ be nonzero and integrable on $X$, and put $F=f/R$. If $\lambda>0$ and $F\leq AB_\lambda$ almost everywhere for some finite $A>0$, then
\[
 s_\lambda=u_\alpha+\lambda^{-1}\log(CV_0/A)
\]
is a genuine exponential subsolution of actual mass $CV_0$. The exponential equation has a greatest genuine solution of finite positive mass. This conclusion holds for every $\lambda>0$ if $F$ is locally essentially bounded off $X_0\times\{o\}$ and
\[
 \lim_{T\to+\infty}\operatorname*{ess\,sup}_{x\in X_0,\ L(y)>T}
 \frac{\log^+(F(x,y)/q_\alpha(y))}{L(y)^\alpha}=0,
 \qquad \log^+t=\log\max(1,t).
\]
\end{theorem}
\begin{proof}
The product measure in \cref{thm:zero-lower-mass-product-equation} is $\mathrm{MA}^{\mathrm{ext}}_H(u_\alpha)=CV_0q_\alpha\mu$. Since $\mathrm e^{\lambda u_\alpha}B_\lambda=q_\alpha$ off the pole,
\[
 \mathrm e^{\lambda s_\lambda}f\,\mathrm dV
 \leq(CV_0/A)\mathrm e^{\lambda u_\alpha}AB_\lambda\mu
 =\mathrm{MA}^{\mathrm{ext}}_H(s_\lambda).
\]
Both measures vanish on the polar fibre. Adding the constant preserves admissibility and mass, and \cref{thm:individual-exponential-subsolution} gives the greatest exact solution. Its mass is finite because its supremum is finite and $f\in L^1$; it is positive because a genuine quasi-plurisubharmonic function is finite almost everywhere and $f$ has positive integral.

Under the limit condition, $F/q_\alpha\leq\mathrm e^{(\lambda/2)L^\alpha}$ almost everywhere sufficiently near the pole. On the remaining compact set the ratio is bounded. Increasing a finite constant $A$ therefore gives $F\leq AB_\lambda$ on all of $X$.
\end{proof}

\begin{theorem}\label{thm:zero-threshold-solvable-concentration}\index{Monge--Ampere density@Monge--Amp\`ere density!solvable concentration}\index{Subsolution!zero threshold exponential equation}
Let $F_j$ be any smooth strictly positive probability densities relative to $\mathrm dV_0=V_0^{-1}h_0^m$ satisfying $c_{h_0}(F_j)\leq2^{-2j}/C$. There are smooth strictly positive probability densities $g_j$ relative to $\kappa^k$ such that
\[
 f_j=RF_jg_j,\qquad f=\sum_{j\geq1}2^{-j}f_j,
 \qquad c_H(f_j)\leq2^{-2j},\qquad b_H(f)=0.
\]
The function $f$ is finite everywhere, smooth off the polar fibre, has integral one and is bounded below by a positive constant. This single density has a genuine exponential subsolution and a greatest finite-mass exact solution for every $\lambda>0$. Every genuine exact solution is unbounded.
\end{theorem}
\begin{proof}
Put $w_j=2^{-j}$ and $M_j=\max F_j$. Using an increasing sequence of smooth coordinate seminorms on $X_0$, choose $0<\varepsilon_j\leq1/2$ so that
\[
 w_j\varepsilon_jM_j\leq b_0 2^{-j-2},\qquad
 w_j\varepsilon_j\|F_j\|_{C^j}\leq2^{-j}.
\]
Then $E=\sum_jw_j\varepsilon_jF_j$ is smooth and satisfies $0<E\leq b_0/4$.

Choose a smooth function $0\leq\rho\leq1$ on $\mathbb R$, supported in $[0,1]$, positive inside and flat at the endpoints. Starting with $A_0=0$, inductively choose $A_j\geq2$, $A_j>A_{j-1}+2$, with
\[
 Z_j=\int_{A_j}^{A_j+1}\rho(L-A_j)
       \mathrm e^{(L^\alpha-1)/j}W(L)\,\mathrm dL
       \geq2w_jM_j.
\]
For each fixed $j$, \cref{lem:projective-exponential-barrier-growth} makes the integrand's weighted radial factor uniformly large on sufficiently distant intervals. Thus such a choice is possible. On that annulus put $b_j=\rho(L-A_j)B_{1/j}/Z_j$, extend it by zero elsewhere and set $g_j=\varepsilon_j+(1-\varepsilon_j)b_j$. Flatness makes $b_j$ smooth across the annulus boundaries; it vanishes near the pole and the infinity hyperplane. The integral of $b_j$ is one, so $g_j$ is a smooth strictly positive probability density. Moreover,
\[
 w_jM_j(1-\varepsilon_j)b_j\leq B_{1/j}/2.
\]

Apply \cref{thm:tw-smooth} on $Y$ to solve $(\kappa+\ddc\phi_j)^k=g_j\kappa^k$ with a smooth positive form; closed integration fixes its scalar as one. On $X_0$ the same theorem gives a smooth potential $v_j$ with $(h_0+\ddc v_j)^m=c_{h_0}(F_j)F_j\,\mathrm dV_0$. The product potential $v_j+\phi_j$ satisfies
\[
 (H+\ddc(v_j+\phi_j))^N
 =C c_{h_0}(F_j)F_jg_j\mu
 =C c_{h_0}(F_j)f_j\,\mathrm dV.
\]
The factor dimensions leave only the displayed binomial term. Scalar uniqueness gives $c_H(f_j)=C c_{h_0}(F_j)\leq2^{-2j}$.

Tonelli's theorem gives $\int_Xf\,\mathrm dV=\sum_jw_j=1$. The annuli are disjoint and escape to the pole, so only finitely many bumps meet any compact set outside the polar fibre. Hence $F=E+\sum_jw_j(1-\varepsilon_j)F_jb_j$ is smooth there. At the fibre the pointwise sum equals $E$, so $f$ is finite everywhere. Also $f\geq w_1f_1\geq\frac12\min_Xf_1>0$. If a positive coefficient $a$ were admitted by the maximal-subsolution characterization in \cref{thm:individual-density-subsolution}, then comparison with each smooth component would give
\[
 aw_j\leq c_H(f_j)\leq2^{-2j},\qquad a\leq2^{-j}
 \quad\text{for every }j.
\]
This is impossible; thus $b_H(f)=0$.

Fix $\lambda>0$ and choose $J$ with $1/j\leq\lambda$ for $j\geq J$. Since $u_\alpha\leq0$, one has $B_{1/j}\leq B_\lambda$ for these indices. Disjointness bounds the sum of their contributions by $B_\lambda/2$, and the baseline is at most $B_\lambda/4$. The remaining finitely many bumps have a finite bound
\[
 K_\lambda=\max_Y\sum_{j<J}w_j(1-\varepsilon_j)M_jb_j,
 \qquad D_\lambda=3/4+K_\lambda/b_0.
\]
Thus $F\leq D_\lambda B_\lambda$. The preceding criterion gives the explicit subsolution $u_\alpha+\lambda^{-1}\log(CV_0/D_\lambda)$ and the greatest exact solution. A bounded exponential subsolution would give a positive constant coefficient for $f$, contradicting its zero threshold. Every genuine exact solution is therefore unbounded.
\end{proof}

The constructed density also satisfies the logarithmic criterion: on the $j$th annulus, $F/q_\alpha\leq1/4+\frac12\mathrm e^{(L^\alpha-1)/j}$, and away from all annuli the ratio is at most $1/4$. Approaching the pole through annuli forces $j\to\infty$, giving the required uniform limit. Since $b_H(f)=0$, \cref{prop:individual-normalizer-lp} also shows that $f$ belongs to no $L^p$ for $p>1$.

For the elliptic threefold, \cref{thm:vanishing-lower-mass} supplies smooth positive forms with masses $\delta_\ell\to0$. Normalize their volume densities by $\delta_\ell$ and choose a subsequence with $\delta_{\ell_j}\leq2^{-2j}/C$. Scalar uniqueness gives the required $F_j$ with $c_{h_0}(F_j)=\delta_{\ell_j}$. Thus the construction applies to this threefold times $\mathbf P^k$ for every $k\geq1$. By \cref{prop:lower-mass-balanced-metric}, $V_0=24\pi$, so the explicit subsolutions have mass $96\pi$ on the fourfold and $240\pi$ on the fivefold. These numerical values concern the subsolutions; the greatest exact solutions have their own finite positive masses.

The projective components are chosen after the base sequence. The theorem does not determine solvability for an arbitrary previously prescribed concentration family or give a necessary criterion at zero threshold.

\Needspace{12\baselineskip}
\subsection{Solvable exponents and collapse at zero threshold}

Fix a nonzero nonnegative $f\in L^1(\mathrm dV)$ on a connected compact Hermitian manifold $(X,h)$. Let $\Lambda_h(f)$ consist of the positive exponents for which
\[
 \mathrm{MA}^{\mathrm{ext}}_h(u)=\mathrm e^{\lambda u}f\,\mathrm dV
\]
has a genuine $h$-plurisubharmonic solution. Existence of an exponential subsolution is equivalent to membership by \cref{thm:individual-exponential-subsolution}.

\Needspace{9\baselineskip}
\begin{proposition}\label{prop:solvable-exponents-upper-interval}\index{Monge--Ampere equation@Monge--Amp\`ere equation!solvable exponents}
If $\lambda_0\in\Lambda_h(f)$, then $[\lambda_0,+\infty)\subset\Lambda_h(f)$. Consequently this set is empty, equals $(a,+\infty)$ for some $a\geq0$, or equals $[a,+\infty)$ for some $a>0$. If $b_h(f)>0$, then $\Lambda_h(f)=(0,+\infty)$.
\end{proposition}
\begin{proof}
Let $u$ solve the equation at $\lambda_0$, and put $s=u-\max(0,\sup_Xu)$. Then $s\leq0$ and
\begin{align*}
 \mathrm{MA}^{\mathrm{ext}}_h(s)
 &=\mathrm e^{\lambda_0u}f\,\mathrm dV\\
 &\geq\mathrm e^{\lambda_0s}f\,\mathrm dV
 \geq\mathrm e^{\lambda s}f\,\mathrm dV
 \qquad(\lambda\geq\lambda_0).
\end{align*}
The greatest-solution theorem supplies an exact solution at each such $\lambda$. Every nonempty upper set of positive real numbers has the stated form, determined by its infimum and whether that infimum is attained. The positive-threshold case follows from \cref{cor:individual-positive-threshold-exponential}.
\end{proof}

\Needspace{15\baselineskip}
\begin{proposition}\label{prop:zero-threshold-exponential-mass-collapse}\index{Normalization scalar!zero threshold collapse}\index{Monge--Ampere mass@Monge--Amp\`ere mass!zero threshold exponential limit}
Suppose $b_h(f)=0$. Let $\lambda_j>0$ decrease to zero, and let $u_j$ be any genuine exact solution at $\lambda_j$. Put
\[
 M_j=\sup_Xu_j,\qquad v_j=u_j-M_j,\qquad
 c_j=\mathrm e^{\lambda_jM_j}.
\]
Then
\begin{gather*}
 c_j\longrightarrow0,\qquad \lambda_jM_j\longrightarrow-\infty,\\
 \|\mathrm{MA}^{\mathrm{ext}}_h(u_j)\|_{\mathrm{TV}}
 =\|\mathrm{MA}^{\mathrm{ext}}_h(v_j)\|_{\mathrm{TV}}
 \longrightarrow0.
\end{gather*}
Moreover $\mathrm e^{\lambda_ju_j}\to0$ uniformly on $X$, with value zero at poles, and $u_j\to-\infty$ uniformly in the sense that $\sup_Xu_j\to-\infty$.
\end{proposition}
\begin{proof}
For $g_j=\mathrm e^{\lambda_jv_j}f$, \cref{lem:normalized-exponential-npp} applied to the fixed finite non-pluripolar measure $f\,\mathrm dV$ gives $g_j\to f$ in $L^1$. Each $g_j$ is nonzero, because $v_j$ is finite almost everywhere. Constant invariance gives
\[
 \mathrm{MA}^{\mathrm{ext}}_h(v_j)=c_jg_j\,\mathrm dV,
 \qquad 0<c_j\leq b_h(g_j)
\]
by \cref{thm:individual-density-subsolution}. Upper semicontinuity, \cref{thm:individual-threshold-semicontinuity}, now gives $\limsup_jc_j\leq b_h(f)=0$. Since $0\leq g_j\leq f$, the total mass is at most $c_j\|f\|_1$ and tends to zero. Also $\lambda_jM_j=\log c_j\to-\infty$ and $0\leq\mathrm e^{\lambda_ju_j}\leq c_j$. For all sufficiently large $j$, $\lambda_j<1$ and $\log c_j<0$, so $M_j=(\log c_j)/\lambda_j\leq\log c_j\to-\infty$.
\end{proof}

For the tailored density of \cref{thm:zero-threshold-solvable-concentration}, solutions exist for every positive exponent, so the collapse applies as $\lambda\downarrow0$ to any choice of genuine exact solutions, including the greatest ones. The subsolutions in \cref{thm:product-zero-threshold-exponential-criterion} retain actual mass $CV_0$ for every exponent, while the exact-solution masses tend to zero. Order of potentials does not give order of their total Monge--Amp\`ere masses on a nonclosed background. Convergence of normalized potentials in capacity and an exact equation for a normalized cluster limit remain separate questions at exponent zero. Positive exponent endpoints are governed instead by the compactness criterion below.

\Needspace{12\baselineskip}
\subsection{Positive endpoints and right limits}

For $\lambda\in\Lambda_h(f)$, write $U_\lambda$ for the greatest solution of \cref{thm:individual-exponential-subsolution}. A bounded truncated-density equation gives a common upper barrier when the exponents stay away from zero. Greatest-subsolution maximality then supplies an order estimate, which controls positive endpoints and right limits.

\Needspace{12\baselineskip}
\begin{lemma}\label{lem:positive-exponent-uniform-upper-barrier}\index{Subsolution!common upper barrier}\index{Monge--Ampere equation@Monge--Amp\`ere equation!positive exponent order}
For every $\tau>0$, there is a nonnegative $t_\tau\in\PSH(X,h)\cap C(X)$ such that every genuine exponential subsolution at an exponent $\lambda\geq\tau$ lies below $t_\tau$. Put $K_\tau=\sup_Xt_\tau$. For $\tau\leq\lambda_1\leq\lambda_2$ in $\Lambda_h(f)$,
\[
 U_{\lambda_2}\geq U_{\lambda_1}
             -K_\tau\frac{\lambda_2-\lambda_1}{\lambda_2}.
\]
Consequently $U_\lambda+K_\tau\log(\lambda/\tau)$ is increasing on $\Lambda_h(f)\cap[\tau,+\infty)$.
\end{lemma}
\begin{proof}
Set $q=\min(f,1)$. This is a bounded nonzero density, so \cref{thm:nguyen-exp} gives a continuous solution $w$ of $\mathrm{MA}_h(w)=\mathrm e^{\tau w}q\,\mathrm dV$. Put $t_\tau=w+\max(0,-\inf_Xw)$. Then $t_\tau\geq0$ and, for $\lambda\geq\tau$,
\[
 \mathrm{MA}_h(t_\tau)=\mathrm e^{\tau w}q\,\mathrm dV
       \leq\mathrm e^{\lambda t_\tau}f\,\mathrm dV.
\]
If $u$ is a genuine subsolution at $\lambda$, then on $\{t_\tau<u-\epsilon\}$,
\[
 \mathrm{MA}_h(t_\tau)
       \leq\mathrm e^{-\lambda\epsilon}
                  \mathrm{MA}^{\mathrm{ext}}_h(u-\epsilon).
\]
The bounded first potential admits a genuine $h/2$-envelope, and $u-\epsilon\preceq t_\tau$. Thus \cref{thm:singular-strict-domination} gives $t_\tau\geq u-\epsilon$. Let $\epsilon\downarrow0$.

For the order estimate, put $\delta=K_\tau(\lambda_2-\lambda_1)/\lambda_2$. Since $U_{\lambda_1}\leq K_\tau$, constant invariance gives
\begin{align*}
 \mathrm{MA}^{\mathrm{ext}}_h(U_{\lambda_1}-\delta)
   &=\mathrm e^{\lambda_1U_{\lambda_1}}f\,\mathrm dV\\
   &\geq\mathrm e^{\lambda_2(U_{\lambda_1}-\delta)}f\,\mathrm dV.
\end{align*}
Greatest-subsolution maximality yields $U_{\lambda_2}\geq U_{\lambda_1}-\delta$. Finally $\log x\geq1-x^{-1}$ for $x\geq1$, which proves the logarithmically corrected order.
\end{proof}

\Needspace{16\baselineskip}
\begin{proposition}\label{prop:positive-exponent-endpoint-attainment}\index{Monge--Ampere equation@Monge--Amp\`ere equation!positive endpoint}
Suppose $\Lambda_h(f)\ne\varnothing$ and $a=\inf\Lambda_h(f)>0$. Then $a\in\Lambda_h(f)$ if and only if there are exponents $\lambda_j>a$ decreasing to $a$ and genuine exact solutions $u_j$ at $\lambda_j$ with $\inf_j\sup_Xu_j> -\infty$.

If $a\notin\Lambda_h(f)$, every sequence of genuine exact solutions at exponents $\lambda_j>a$ decreasing to $a$ satisfies
\begin{gather*}
 \sup_Xu_j\longrightarrow-\infty,\qquad
 \mathrm e^{\lambda_ju_j}\longrightarrow0\ \text{uniformly on }X,\\
 \|\mathrm{MA}^{\mathrm{ext}}_h(u_j)\|_{\mathrm{TV}}
                \longrightarrow0.
\end{gather*}
The exponential has value zero at poles.
\end{proposition}
\begin{proof}
If $u_a$ is an exact solution at $a$, put $s=u_a-\max(0,\sup_Xu_a)$. Then $s\leq0$ and $\mathrm{MA}^{\mathrm{ext}}_h(s)\geq\mathrm e^{\lambda s}f\,\mathrm dV$ for every $\lambda\geq a$. Thus $U_\lambda\geq s$, which supplies the required sequence.

Conversely, the lemma with $\tau=a$ gives a common upper bound $K_a$ for $u_j$. Write $M_j=\sup_Xu_j$. Extract a subsequence for which $M_j\to M\in\mathbb R$ and, by \cref{thm:normalized-compactness}, $u_j-M_j\to v\in\PSH(X,h)$ in $L^1$. Then $u_j\to u=v+M$ in $L^1$, with a genuine limit and a common upper bound. After a further extraction the convergence holds almost everywhere. Since $\lambda_j\leq\lambda_1$ and $K_a\geq0$, dominated convergence gives
\[
 \mathrm e^{\lambda_ju_j}f\longrightarrow\mathrm e^{au}f
              \quad\text{in }L^1(\mathrm dV).
\]
Indeed all these densities are at most $\mathrm e^{\lambda_1K_a}f$. The lower-density closure \cref{lem:density-subsolution-lone-closure} gives $\mathrm{MA}^{\mathrm{ext}}_h(u)\geq\mathrm e^{au}f\,\mathrm dV$. The greatest-solution theorem now supplies an exact solution at $a$.

If $a$ is not admitted and the suprema of an exact solution sequence failed to tend to $-\infty$, a subsequence would satisfy the lower-bound condition just proved. Thus $M_j\to-\infty$. Eventually $M_j<0$; since $\lambda_j\geq a>0$,
\[
 0\leq\mathrm e^{\lambda_ju_j}\leq\mathrm e^{aM_j}
                 \longrightarrow0.
\]
The exact mass is at most $\mathrm e^{aM_j}\|f\|_1$, proving the last assertion.
\end{proof}

\Needspace{10\baselineskip}
\begin{proposition}\label{prop:greatest-exponential-right-continuity}\index{Monge--Ampere equation@Monge--Amp\`ere equation!right continuity}\index{Convergence!greatest exponential solutions}
For every $a\in\Lambda_h(f)$, the greatest solutions satisfy
\[
 U_\lambda\longrightarrow U_a
       \quad\text{in }L^1\text{ and in capacity}\quad(\lambda\downarrow a).
\]
Their Monge--Amp\`ere measures converge in total variation to $\mathrm e^{aU_a}f\,\mathrm dV$.
\end{proposition}
\begin{proof}
Put $K=K_a$ from the lemma. For any $\lambda_j>a$ decreasing to $a$, the potentials
\[
 V_j=U_{\lambda_j}+K\log(\lambda_j/a)
\]
decrease and lie above $U_a$. Their limit $V$ is genuine. Dominated convergence between the integrable bounds $U_a$ and $V_1$ gives $L^1$ convergence, and \cref{lem:monotone-psh-capacity} gives capacity convergence. Since the added constants tend to zero, $U_{\lambda_j}\to V$ in both topologies.

Apply \cref{thm:exponential-closedness} to the nonpositive potentials $U_{\lambda_j}-K$, the exponents $\lambda_j\to a>0$ and the densities $\mathrm e^{\lambda_jK}f\to\mathrm e^{aK}f$ in $L^1$. It gives $\mathrm{MA}^{\mathrm{ext}}_h(V)=\mathrm e^{aV}f\,\mathrm dV$, with total-variation convergence of the equation measures. Hence $V\leq U_a$ by greatest-solution maximality, while the lower bound gives $V\geq U_a$. Thus $V=U_a$. The argument applies to every decreasing sequence of exponents, giving the stated right limit.
\end{proof}

The endpoint test is expressed in terms of existing solutions. An intrinsic criterion on $f$, the existence or nonexistence of a positive excluded endpoint, and continuity from the left remain separate questions. The right limit selects greatest solutions; it does not identify limits of arbitrary solutions with other singularities.
